% Corpus travail-emploi : source, version et structure (ticket #27, domaine D1).
% Compilation depuis docs/corpus : latexmk -xelatex CORPUS.tex
\documentclass[11pt,a4paper]{article}

% ---------------------------------------------------------------- langue & fontes
\usepackage{fontspec}
\usepackage[french]{babel}
\frenchsetup{SmallCapsFigTabCaptions=false,ThinSpaceInFrenchNumbers=true}
\setmainfont{LibertinusSerif-Regular.otf}[
  BoldFont=LibertinusSerif-Bold.otf, ItalicFont=LibertinusSerif-Italic.otf,
  BoldItalicFont=LibertinusSerif-BoldItalic.otf, Numbers=OldStyle]
\setsansfont{Inter-Regular.otf}[
  BoldFont=Inter-Bold.otf, ItalicFont=Inter-Italic.otf,
  BoldItalicFont=Inter-BoldItalic.otf, Scale=0.94]
\setmonofont{FiraMono-Regular.otf}[BoldFont=FiraMono-Bold.otf, Scale=0.84]
\newfontfamily\tabfont{Inter-Regular.otf}[BoldFont=Inter-Bold.otf,
  ItalicFont=Inter-Italic.otf, BoldItalicFont=Inter-BoldItalic.otf,
  Scale=0.94, Numbers=Lining]

% ---------------------------------------------------------------- mise en page
\usepackage[a4paper,top=2.5cm,bottom=2.3cm,left=2.1cm,right=2.1cm,
            headheight=17pt,headsep=13pt,footskip=20pt]{geometry}
\usepackage{microtype}
\usepackage{xcolor}
\usepackage{booktabs,tabularx,longtable,array,colortbl,multirow}
\usepackage{titlesec,titletoc}
\usepackage{fancyhdr}
\usepackage{enumitem}
\usepackage[most]{tcolorbox}
\usepackage{fancyvrb}
\usepackage{pifont}
\usepackage{fontawesome5}
\usepackage{needspace}
\usepackage{ragged2e}
\usepackage[hidelinks]{hyperref}

% ---------------------------------------------------------------- palette
\definecolor{encre}   {HTML}{16181D}
\definecolor{accent}  {HTML}{1F4E5F}
\definecolor{accent2} {HTML}{8C4A2F}
\definecolor{gris}    {HTML}{6B7280}
\definecolor{filet}   {HTML}{D5DADE}
\definecolor{fondclair}{HTML}{F3F6F7}
\definecolor{fondcode}{HTML}{F7F7F4}
\definecolor{vert}    {HTML}{17693F}
\definecolor{rouge}   {HTML}{9B2C2C}
\definecolor{ambre}   {HTML}{A86A12}
\color{encre}

\hypersetup{colorlinks=true,linkcolor=accent,urlcolor=accent,
  pdftitle={Reunion de cadrage - Projet NLP 2},
  pdfsubject={Jalon J2 - choix des cas d'usage, cadrage, outils novateurs},
  pdfauthor={Equipe projet NLP 2 - M2 Data et IA, FGES}}

% ---------------------------------------------------------------- symboles
\newcommand{\okmark}   {\textcolor{vert}{\ding{51}}}
\newcommand{\komark}   {\textcolor{rouge}{\ding{55}}}
\newcommand{\warnmark} {\textcolor{ambre}{\faExclamationTriangle}}
\newcommand{\checkbox} {\textcolor{gris}{\ding{113}}}
\newcommand{\cutmark}  {\textcolor{accent2}{\ding{34}}}
\newcommand{\starmark} {\textcolor{ambre}{\ding{72}}}
\newcommand{\wipmark}  {\textcolor{ambre}{\faTools}}
\newcommand{\code}[1]  {{\ttfamily #1}}
\newcommand{\rulebreak}{}

% ---------------------------------------------------------------- titres
\titleformat{\section}[hang]
  {\sffamily\bfseries\LARGE\color{accent}}{\thesection}{13pt}{}
  [\vspace{2pt}{\color{filet}\titlerule[0.9pt]}]
\titlespacing*{\section}{0pt}{24pt plus 4pt minus 2pt}{11pt}
\titleformat{\subsection}
  {\sffamily\bfseries\large\color{encre}}{\thesubsection}{9pt}{}
\titlespacing*{\subsection}{0pt}{16pt plus 3pt minus 1pt}{6pt}
\titleformat{\subsubsection}
  {\sffamily\bfseries\normalsize\color{accent2}}{}{0pt}{}
\titlespacing*{\subsubsection}{0pt}{12pt plus 2pt minus 1pt}{4pt}
\setcounter{secnumdepth}{2}
\setcounter{tocdepth}{2}

% ---------------------------------------------------------------- en-têtes
\pagestyle{fancy}\fancyhf{}
% en-tête courant, redéfini par build.py selon le document
\newcommand{\entetecourant}{Projet NLP 2}
\fancyhead[L]{\sffamily\scriptsize\color{gris}\entetecourant}
\fancyhead[R]{\sffamily\scriptsize\color{gris}\thepage}
\renewcommand{\headrulewidth}{0pt}
\renewcommand{\headrule}{\vspace{-4pt}{\color{filet}\rule{\headwidth}{0.4pt}}}
\fancypagestyle{plain}{\fancyhf{}\renewcommand{\headrule}{}%
  \fancyfoot[C]{\sffamily\scriptsize\color{gris}\thepage}}

% ---------------------------------------------------------------- paragraphes & listes
\setlength{\parindent}{0pt}
\setlength{\parskip}{5.5pt plus 1.5pt minus 0.5pt}
\linespread{1.05}
\raggedbottom
\widowpenalty=10000 \clubpenalty=10000
\setlist[itemize]{leftmargin=1.15em,itemsep=2pt,topsep=3pt,parsep=0pt,label=\textcolor{accent}{\textbullet}}
\setlist[enumerate]{leftmargin=1.5em,itemsep=2.5pt,topsep=3pt,parsep=0pt,
  label=\textcolor{accent}{\sffamily\bfseries\arabic*.}}

% ---------------------------------------------------------------- tableaux
% \hspace{0pt} : glue initiale, sans quoi TeX refuse de couper le premier
% mot de la cellule — cause de débordements en colonne étroite.
\newcolumntype{L}[1]{>{\raggedright\arraybackslash
  \rightskip=0pt plus 1fil\relax\hspace{0pt}}p{#1}}
% \hspace{0pt} : sans cette glue, le whatsit de \color empêche TeX
% de couper le premier mot de la cellule (débordements en en-tête).
\newcommand{\thead}[1]{{\tabfont\bfseries\color{accent}\hspace{0pt}#1}}
\newenvironment{tabletype}
  {\par\addvspace{7pt}\begingroup
   \tabfont\setlength{\tabcolsep}{4pt}\renewcommand{\arraystretch}{1.26}
   \setlength{\LTleft}{0pt}\setlength{\LTright}{0pt}
   \setlength{\extrarowheight}{1pt}
   % élasticité infinie à droite : TeX ne coupe plus un mot pour « remplir »,
   % mais coupe encore lorsque c'est le seul moyen d'éviter un débordement.
   % césure interdite : les largeurs minimales sont mesurées, chaque mot tient.
   % Un débordement, s'il en restait un, serait signalé par le journal.
   \hyphenpenalty=9000\exhyphenpenalty=9000
   \emergencystretch=0pt\hbadness=10000\relax}
  {\endgroup\par\addvspace{9pt}}
\setlength{\aboverulesep}{0.5pt}\setlength{\belowrulesep}{0.6pt}
\arrayrulecolor{filet}
\setlength{\heavyrulewidth}{0.9pt}\setlength{\lightrulewidth}{0.45pt}

% ---------------------------------------------------------------- encadrés
\newtcolorbox{callout}{enhanced,breakable,
  colback=fondclair,colframe=accent,boxrule=0pt,leftrule=2.6pt,
  sharp corners,left=11pt,right=11pt,top=8pt,bottom=8pt,
  before skip=10pt,after skip=11pt,parbox=false}
\newtcolorbox{codebox}{enhanced,breakable,
  colback=fondcode,colframe=filet,boxrule=0.5pt,arc=2pt,
  left=9pt,right=9pt,top=7pt,bottom=7pt,
  before skip=10pt,after skip=11pt}

% ---------------------------------------------------------------- sommaire
\renewcommand{\contentsname}{Sommaire}
\titlecontents{section}[0pt]
  {\addvspace{4pt}\sffamily\bfseries\color{encre}}
  {\contentslabel[\thecontentslabel]{1.6em}}{}
  {\titlerule*[0.7pc]{\color{filet}.}\contentspage}
\titlecontents{subsection}[1.9em]
  {\sffamily\small\color{gris}}
  {\contentslabel[\thecontentslabel]{2.3em}}{}
  {\titlerule*[0.7pc]{\color{filet}.}\contentspage}

% ---------------------------------------------------------------- divers
\newlength{\metalab}
\newlength{\tw}\setlength{\tw}{\textwidth}
\newcommand{\settw}[1]{\setlength{\tw}{\dimexpr\textwidth-#1pt\relax}}
% tolérance relâchée : sans elle, TeX préfère déborder de quelques points
% plutôt que de desserrer une ligne contenant un long jeton en chasse fixe.
\emergencystretch=1.6em

\usepackage{tikz}
\usetikzlibrary{arrows.meta,positioning,fit,backgrounds,calc}

\hypersetup{pdftitle={Corpus travail-emploi - source, version et structure},
  pdfsubject={Ticket 27 - domaine D1, corpus et segmentation},
  pdfauthor={Equipe Data Tigers - M2 Data et IA, FGES}}
\renewcommand{\entetecourant}{Corpus travail-emploi \textbullet{} Équipe Data Tigers}

% ---------------------------------------------------------------- figures
\tikzset{
  sf/.style={font=\sffamily\fontsize{7.6}{9.2}\selectfont, text=encre},
  fl/.style={-{Stealth[length=4.2pt,width=3.4pt]}, line width=0.65pt, draw=encre!80},
  etape/.style={sf, rectangle, rounded corners=2pt, draw=accent, fill=white,
    line width=0.6pt, align=center, text width=3.0cm, minimum height=1.35cm, inner sep=4pt},
  bandeau/.style={sf, rectangle, fill=accent, text=white, anchor=north west,
    font=\sffamily\bfseries\fontsize{7.8}{9.4}\selectfont, inner xsep=6pt, inner ysep=4pt},
  section/.style={rectangle, rounded corners=2pt, draw=filet, fill=fondclair, line width=0.5pt},
  pastille/.style={sf, rectangle, rounded corners=1.5pt, draw=accent, fill=white,
    minimum width=0.44cm, minimum height=0.44cm, inner sep=0pt,
    font=\sffamily\bfseries\fontsize{7}{8}\selectfont},
  choisie/.style={pastille, fill=accent2, draw=accent2, text=white},
  morceau/.style={sf, rectangle, rounded corners=2pt, line width=0.6pt, align=left,
    anchor=north west, inner sep=5pt},
}

\newcolumntype{Y}{>{\raggedright\arraybackslash\rightskip=0pt plus 1fil\relax\hspace{0pt}}X}
\newcommand{\var}[1]{{\ttfamily\bfseries\color{accent2}#1}}
\newcommand{\ex}[1]{{\color{encre}#1}}

\begin{document}
\thispagestyle{empty}
\begingroup\sffamily
{\color{accent}\rule{\textwidth}{2.6pt}}\par\vspace{13pt}
{\fontsize{27}{31}\selectfont\bfseries\color{encre} Corpus travail-emploi}\par\vspace{7pt}
{\fontsize{13.5}{17}\selectfont\color{gris} Équipe Data Tigers \textbullet{} NLP~2 \textbullet{} M2 Data \& IA \textbullet{} FGES}\par\vspace{5pt}
{\fontsize{11}{14}\selectfont\color{accent2}\bfseries Source, version et structure des données}\par\vspace{12pt}
{\color{filet}\rule{\textwidth}{0.6pt}}\par\vspace{14pt}
\endgroup

\begingroup\setlength{\parskip}{0pt}
\settowidth{\metalab}{\tabfont\bfseries Responsables}
\begin{tabular}{@{}>{\tabfont\bfseries\color{accent}}l@{\hspace{16pt}}p{\dimexpr\textwidth-\metalab-18pt\relax}@{}}
Ticket & \#27, télécharger, filtrer et contrôler le corpus \\[3.5pt]
Domaine & D1, corpus et segmentation \\[3.5pt]
Responsables & Maïmouna SIGNATE et Vaneck DAGAR, soutien Jibril BENSALEM \\[3.5pt]
Statut & document de travail, tâches 2 à 8 du ticket \\[3.5pt]
\end{tabular}\par\vspace{16pt}\endgroup

\begingroup\small\color{gris}
Ce document couvre les tâches 2 et 3 du ticket : documenter la source, la version et la méthode
d'obtention du corpus, puis identifier son format et sa structure. La section~\ref{sec:filtrage}
décrit le filtrage, le nombre de passages conservés et la volumétrie après filtrage (tâches 4 à 6). La
section~\ref{sec:controle} présente le contrôle manuel d'un échantillon (tâche 7), et la
section~\ref{sec:anomalies} recense les anomalies et les limites observées (tâche 8). Sauf mention
contraire, tous les chiffres proviennent d'une mesure effectuée le 26 septembre 2026 sur le fichier
décrit en section~\ref{sec:fichier}.
\par\endgroup
\vspace{10pt}{\color{filet}\rule{\textwidth}{0.4pt}}\par\vspace{6pt}

% ====================================================================
\section{Source, version et méthode d'obtention}

\subsection{Provenance}

Le corpus retenu au jalon J2 est le jeu de données public \code{AgentPublic/travail-emploi}, publié
sur Hugging Face par l'administration française. Son contenu provient d'un seul site source, celui
du ministère du Travail. Il nous parvient au terme d'une chaîne de trois maillons, représentée en
figure~\ref{fig:provenance}.

\begin{figure}[htbp]
\centering
\begin{tikzpicture}
  \node[etape] (site) at (0,0)
    {{\bfseries Site du ministère}\\[2pt]{\color{gris}travail-emploi.gouv.fr\\fiches pratiques}};
  \node[etape] (sg) at (4.4,0)
    {{\bfseries SocialGouv}\\[2pt]{\color{gris}dépôt GitHub\\\ttfamily fiches-travail-data}};
  \node[etape] (mt) at (8.8,0)
    {{\bfseries MediaTech (Etalab)}\\[2pt]{\color{gris}découpage en passages\\et vectorisation}};
  \node[etape, fill=accent!9] (hf) at (13.2,0)
    {{\bfseries Hugging Face}\\[2pt]{\color{gris}\ttfamily AgentPublic/\\travail-emploi}};
  \draw[fl] (site) -- (sg);
  \draw[fl] (sg) -- (mt);
  \draw[fl] (mt) -- (hf);
\end{tikzpicture}
\caption{Chaîne de provenance du corpus. Nous utilisons le dernier maillon, sans retraitement à ce stade.}
\label{fig:provenance}
\end{figure}

\begin{tabletype}
\begin{tabularx}{\textwidth}{@{}L{4cm}Y@{}}
\toprule
\thead{Élément} & \thead{Valeur} \\
\midrule
Jeu de données & \code{AgentPublic/travail-emploi} \\
Adresse & \url{https://huggingface.co/datasets/AgentPublic/travail-emploi} \\
Éditeur & AgentPublic, collection MediaTech (Etalab) \\
Licence & Etalab 2.0, licence ouverte \\
Contenu d'origine & fiches pratiques de \url{https://travail-emploi.gouv.fr/} \\
Code de production & \url{https://github.com/etalab-ia/mediatech} \\
\bottomrule
\end{tabularx}
\end{tabletype}

\subsection{Versions disponibles et version retenue}

Le dépôt Hugging Face ne publie pas un fichier unique mais une série d'instantanés datés, un par mise
à jour, rangés sous \code{data/travail-emploi-AAAAMMJJ/}. Au 26 septembre 2026, il en contient 28, du
24 octobre 2025 au 11 septembre 2026, plus un dossier \code{data/travail-emploi-latest/} qui suit la
dernière version et change à chaque mise à jour.

\textbf{Version retenue : l'instantané du 11 septembre 2026.} C'est la version mesurée au jalon J2,
dont elle reproduit exactement la volumétrie de 5\,702 passages. Un instantané daté ne change plus,
contrairement au dossier \code{latest}, ce qui garantit que tous les domaines travaillent sur les mêmes
données.

L'historique du dépôt confirme la correspondance : jusqu'au 25 septembre 2026, le dossier \code{latest}
contenait cet instantané. Il a été renommé \code{travail-emploi-20260911} le 25 septembre, au moment de la
publication d'une nouvelle version.

\begin{tabletype}
\begin{tabularx}{\textwidth}{@{}L{5.6cm}L{2cm}L{2cm}Y@{}}
\toprule
\thead{Version} & \thead{Passages} & \thead{Documents} & \thead{Écart} \\
\midrule
\code{travail-emploi-20260911}, retenue & 5\,702 & 261 & référence \\
\code{travail-emploi-latest} au 26/09/2026 & 5\,703 & 261 & 1 passage ajouté, 57 modifiés \\
\bottomrule
\end{tabularx}
\end{tabletype}

\subsection{Fichier de référence}\label{sec:fichier}

\begin{tabletype}
\begin{tabularx}{\textwidth}{@{}L{4cm}Y@{}}
\toprule
\thead{Élément} & \thead{Valeur} \\
\midrule
Chemin sur Hugging Face & \code{data/travail-emploi-20260911/travail\_emploi\_part\_0.parquet} \\
Commit figé & \code{98744a63d40d982ff8b5345ca635ad543a412163}, du 25/09/2026 \\
Taille & 31\,031\,404 octets, soit 29,6 Mio \\
Empreinte SHA-256 & \code{b76ab4c4408b737fadaade549dfa7734}\newline\code{ec9b3343b7c19f0285cf3c6cd858ef4f} \\
Date d'obtention & 26 septembre 2026 \\
\bottomrule
\end{tabularx}
\end{tabletype}

\subsection{Méthode d'obtention}

Le fichier est téléchargé directement par HTTPS, en désignant le commit plutôt que la branche
\code{main}. L'adresse reste ainsi valable même si le dépôt évolue. Les commandes suivantes, lancées
depuis la racine du dépôt, reconstruisent le fichier à l'identique.

\begin{codebox}
\begin{Verbatim}[fontsize=\footnotesize]
mkdir -p data/raw/travail-emploi-20260911
curl -L -o data/raw/travail-emploi-20260911/travail_emploi_part_0.parquet \
  "https://huggingface.co/datasets/AgentPublic/travail-emploi/resolve/\
98744a63d40d982ff8b5345ca635ad543a412163/\
data/travail-emploi-20260911/travail_emploi_part_0.parquet"
sha256sum data/raw/travail-emploi-20260911/travail_emploi_part_0.parquet
\end{Verbatim}
\end{codebox}

L'empreinte obtenue doit être identique à celle de la section~\ref{sec:fichier}. Une empreinte
différente signale un fichier corrompu ou une autre version.

Le dossier \code{data/} est exclu du dépôt Git par le \code{.gitignore} : le fichier pèse 29,6 Mio et se
reconstruit par les commandes ci-dessus. Aucune bibliothèque spécifique n'est requise pour le
téléchargement. La lecture utilise \code{pandas} et \code{pyarrow}.

\textbf{Reproduction complète en une commande.} Le script \code{scripts/prepare\_corpus.py} enchaîne
toutes les étapes : il télécharge le fichier s'il est absent, vérifie son empreinte, applique les
règles de filtrage (section~\ref{sec:filtrage}), compte les tokens de chaque passage et écrit ses
résultats. Lancé depuis la racine du dépôt, il produit toujours le même corpus. Il requiert les
bibliothèques \code{pandas}, \code{pyarrow} et \code{tokenizers} ; le script de tirage de l'échantillon
de contrôle (section~\ref{sec:controle}) requiert en outre \code{openpyxl}.

\begin{codebox}
\begin{Verbatim}[fontsize=\footnotesize]
python scripts/prepare_corpus.py                  # doublons de R4 marqués
python scripts/prepare_corpus.py --sans-doublons  # doublons de R4 retirés
\end{Verbatim}
\end{codebox}

Il écrit trois fichiers dans \code{data/processed/} :

\begin{itemize}
\item \code{travail\_emploi\_filtre.parquet} : le corpus filtré ;
\item \code{passages\_retires.csv} : la liste des passages retirés, avec la règle appliquée ;
\item \code{stats\_filtrage.json} : les mesures avant et après filtrage.
\end{itemize}

% ====================================================================
\section{Format et structure des données}

\subsection{Format}

\begin{tabletype}
\begin{tabularx}{\textwidth}{@{}L{4cm}Y@{}}
\toprule
\thead{Élément} & \thead{Valeur} \\
\midrule
Format & Apache Parquet, un seul fichier \\
Compression & ZSTD \\
Organisation & un seul groupe de lignes, trié par \code{doc\_id} puis \code{chunk\_index} \\
Granularité & une ligne par passage (\emph{chunk}) \\
Volumétrie & 5\,702 passages issus de 261 documents, 14 colonnes \\
\bottomrule
\end{tabularx}
\end{tabletype}

\subsection{Organisation des textes}

Le corpus s'organise en trois niveaux emboîtés, représentés en figure~\ref{fig:hierarchie} sur un
exemple réel, la fiche « Le bulletin de paie ».

\begin{description}[leftmargin=0pt, itemsep=3pt, font=\sffamily\bfseries\color{accent}]
\item[Fiche.] Une fiche est une page du site travail-emploi.gouv.fr qui explique un sujet précis du
droit du travail : le CDD, le bulletin de paie, le burn-out. Le corpus en compte 261. Chaque fiche est un
\emph{document}, identifié par \code{doc\_id}.
\item[Section.] Une fiche est organisée en sections, qui sont ses intertitres, par exemple « Et les
mentions interdites ? ». L'intitulé de la section figure dans la colonne \code{context}. Le corpus compte
2\,115 sections.
\item[Passage.] Chaque section est découpée en un ou plusieurs passages (\emph{chunks}) selon sa
longueur. Un passage correspond à une ligne du fichier, et son contenu se trouve dans la colonne
\code{text}. Un passage n'appartient jamais à deux sections.
\end{description}

\begin{figure}[htbp]
\centering
\begin{tikzpicture}
  % ------------------------------------------------ site
  \begin{scope}[on background layer]
    \node[rectangle, rounded corners=4pt, fill=fondclair, draw=filet, line width=0.5pt,
          fit={(0,0.35) (16.8,-12.15)}, inner sep=0pt] {};
  \end{scope}
  \node[sf, anchor=north west, font=\sffamily\bfseries\fontsize{8.4}{10}\selectfont, text=accent]
    at (0.25,0.25) {SITE SOURCE \quad travail-emploi.gouv.fr};
  \node[sf, anchor=north east, text=gris] at (16.55,0.25) {un seul site, 261 fiches};

  % ------------------------------------------------ fiche
  \draw[accent, line width=0.8pt, fill=white, rounded corners=3pt] (0.3,-0.45) rectangle (9.95,-11.85);
  \node[bandeau, minimum width=9.65cm] at (0.3,-0.45)
    {FICHE = une page du site \hfill identifiée par \texttt{doc\_id}};
  \foreach \y/\nom/\role/\exemple in {
      -1.5/doc\_id/identifiant de la fiche/{\ttfamily article374540},
      -2.05/title/titre de la page/{Le bulletin de paie},
      -2.65/url/adresse de la page/{\ttfamily\fontsize{6.8}{8}\selectfont travail-emploi.gouv.fr/\\le-bulletin-de-paie},
      -3.25/date/dernière mise à jour/{02/03/2026},
      -3.9/introduction/{paragraphe d'en-tête,\\résumé de la fiche}/{« Au moment du versement de son salaire, un bulletin de paie doit être remis… »}} {
    \node[sf, anchor=west, inner sep=0pt] at (0.55,\y) {\var{\nom}};
    \node[sf, anchor=west, inner sep=0pt, align=left, text=gris] at (2.55,\y) {\role};
    \node[sf, anchor=west, inner sep=0pt, align=left, text width=4.25cm] at (5.45,\y) {\exemple};
  }
  \node[sf, anchor=north west, text=gris, text width=9.2cm, font=\sffamily\itshape\fontsize{7}{8.4}\selectfont]
    at (0.55,-4.55) {Ces champs décrivent la fiche et sont recopiés sur chacun de ses passages.};

  \node[sf, anchor=north west, font=\sffamily\bfseries\fontsize{7.6}{9}\selectfont, text=accent]
    at (0.55,-5.1) {SECTIONS \quad\mdseries\color{gris}colonne \texttt{context}};
  \node[sf, anchor=north west, font=\sffamily\bfseries\fontsize{7.6}{9}\selectfont, text=accent]
    at (6.1,-5.1) {PASSAGES \quad\mdseries\color{gris}\texttt{chunk\_index}};

  % bandes de sections : y centre, intitulé, liste des passages
  \foreach \y/\titre/\liste in {
      -6.0/{« Le bulletin de paie »}/{1},
      -6.9/{« Quelles sont les mentions obligatoires ? »}/{2,3,4,5,6,7},
      -7.8/{« Documents annexés au bulletin de paie »}/{8,9},
      -8.7/{« Et les mentions interdites ? »}/{10},
      -10.5/{« À lire sur service-public.gouv.fr »}/{27}} {
    \draw[section] (0.55,\y+0.36) rectangle (9.75,\y-0.36);
    \node[sf, anchor=west, text width=5.1cm, font=\sffamily\itshape\fontsize{7.4}{8.8}\selectfont] at (0.7,\y) {\titre};
    \foreach \n [count=\i] in \liste {
      \node[pastille] at ({6.35+0.55*(\i-1)},\y) {\n};
    }
  }
  \draw[section, dashed] (0.55,-9.6+0.36) rectangle (9.75,-9.6-0.36);
  \node[sf, anchor=west, text=gris, font=\sffamily\itshape\fontsize{7.4}{8.8}\selectfont] at (0.7,-9.6)
    {8 autres sections, passages 11 à 26};
  \node[choisie] (p10) at (6.35,-8.7) {10};

  \node[sf, anchor=north west, text=gris, text width=9.2cm, font=\sffamily\itshape\fontsize{7}{8.4}\selectfont]
    at (0.55,-11.05) {Cette fiche compte 13 sections et 27 passages. Une section courte tient en un passage,
    une section longue est coupée en plusieurs.};

  % ------------------------------------------------ passage
  \draw[accent2, line width=0.8pt, fill=white, rounded corners=3pt] (10.45,-0.45) rectangle (16.5,-11.85);
  \node[bandeau, fill=accent2, minimum width=6.05cm] at (10.45,-0.45) {PASSAGE n° 10 = une ligne du fichier};
  \node[sf, anchor=north west, text width=5.6cm, inner sep=0pt] at (10.7,-1.2) {%
    \setlength{\parskip}{0pt}%
    \var{chunk\_id}\\ {\color{gris}identifiant unique du passage}\\
    \ex{\ttfamily article374540\_10}\\ {\color{gris}\fontsize{6.8}{8}\selectfont = \texttt{doc\_id} + « \_ » + \texttt{chunk\_index}}\\[6pt]
    \var{chunk\_index}\\ {\color{gris}position du passage dans la fiche}\\ \ex{10}\\[6pt]
    \var{context}\\ {\color{gris}section d'où vient le passage}\\ \ex{« Et les mentions interdites ? »}\\[6pt]
    \var{text}\\ {\color{gris}contenu du passage, 304 caractères}\\
    \ex{« Aucune mention relative à l'exercice du droit de grève et à l'activité de représentation des salariés ne doit figurer… »}\\[6pt]
    \var{chunk\_text}\\ {\color{gris}\texttt{title} + \texttt{introduction} + \texttt{text}, 619 caractères}\\[6pt]
    \var{chunk\_xxh64}\\ {\color{gris}empreinte de \texttt{chunk\_text}}\\ \ex{\ttfamily 3e31bce72789afd2}\\[6pt]
    \var{embeddings\_bge-m3}\\ {\color{gris}vecteur de 1\,024 nombres calculé par bge-m3}\\ \ex{\ttfamily [-0.0434, 0.0103, …]}\\[8pt]
    {\color{gris}\fontsize{6.8}{8}\selectfont Deux colonnes sont constantes : \texttt{surtitle} vaut « Travail-Emploi » et \texttt{source} vaut « travail-emploi ».}};

  \draw[accent2, line width=0.9pt, dashed, -{Stealth[length=4.5pt,width=3.6pt]}] (p10.east) -- (10.45,-8.7);
\end{tikzpicture}
\caption{Organisation des textes et rôle de chaque variable, sur l'exemple réel de la fiche
« Le bulletin de paie ». À gauche, la fiche et ses sections ; chaque pastille est un passage, numéroté
par \code{chunk\_index}. À droite, le détail du passage n°~10.}
\label{fig:hierarchie}
\end{figure}

\subsection{Dictionnaire des variables}

Le fichier compte 14 colonnes. Le tableau les regroupe par niveau : ce que la colonne décrit est
indiqué par son niveau, fiche ou passage. Les longueurs sont en caractères.

\begin{tabletype}
\begin{tabularx}{\textwidth}{@{}L{3.35cm}L{1.3cm}L{1.6cm}L{2.5cm}Y@{}}
\toprule
\thead{Colonne} & \thead{Type} & \thead{Distinctes} & \thead{Long. min/moy/max} & \thead{Signification} \\
\midrule
\multicolumn{5}{@{}l}{\textcolor{accent}{\bfseries Niveau fiche, recopié sur chaque passage}} \\
\code{doc\_id} & texte & 261 & 5 / 13 / 14 & identifiant de la fiche sur le site source \\
\code{title} & texte & 261 & 8 / 47 / 136 & titre de la fiche \\
\code{url} & texte & 261 & 39 / 77 / 160 & adresse de la fiche sur travail-emploi.gouv.fr \\
\code{date} & texte & 140 & 10 & publication ou mise à jour, format \code{JJ/MM/AAAA} \\
\code{introduction} & texte & 255 & 0 / 572 / 2\,056 & paragraphe d'en-tête de la fiche \\
\midrule
\multicolumn{5}{@{}l}{\textcolor{accent}{\bfseries Niveau passage}} \\
\code{chunk\_id} & texte & 5\,702 & 7 / 16 / 17 & identifiant unique du passage \\
\code{chunk\_index} & entier & 199 & de 1 à 199 & rang du passage dans sa fiche, à partir de 1 \\
\code{context} & liste & 1\,267 & 10 / 54 / 267 & intitulé de la section, liste d'un seul élément \\
\code{text} & texte & 5\,503 & 4 / 929 / 1\,464 & contenu du passage \\
\code{chunk\_text} & texte & 5\,702 & 140 / 1\,549 / 3\,438 & titre, introduction et contenu réunis (figure~\ref{fig:composition}) \\
\code{chunk\_xxh64} & texte & 5\,702 & 16 & empreinte XXH64 de \code{chunk\_text} \\
\code{embeddings\_bge-m3} & texte & 5\,702 & — & vecteur de 1\,024 dimensions, chaîne JSON \\
\midrule
\multicolumn{5}{@{}l}{\textcolor{accent}{\bfseries Constantes}} \\
\code{surtitle} & texte & 1 & — & toujours « Travail-Emploi » \\
\code{source} & texte & 1 & — & toujours « travail-emploi » \\
\bottomrule
\end{tabularx}
\end{tabletype}

\subsection{Relations entre les variables}

\textbf{Fiche et passages.} Chaque fiche a un titre unique et une URL unique. Ses passages sont
numérotés de 1 à $n$ sans trou par \code{chunk\_index}. Une fiche compte entre 1 et 199 passages, 15 en
médiane.

\textbf{Identifiant du passage.} \code{chunk\_id} est toujours formé de \code{doc\_id}, d'un tiret bas et de
\code{chunk\_index}. \code{chunk\_index} seul n'identifie rien : le numéro 10 existe dans de nombreuses
fiches.

\textbf{Sections et passages.} Les 5\,702 passages se répartissent en 2\,115 sections : 1\,229 tiennent
en un seul passage, 710 en deux à cinq, et 176 en plus de cinq, jusqu'à 67 pour une même section.

\textbf{Texte brut et texte composé.} \code{text} contient le seul contenu du passage.
\code{chunk\_text} en est une version enrichie, composée par MediaTech pour l'encodage : sur trois lignes,
le titre de la fiche, son introduction, puis le contenu du passage (figure~\ref{fig:composition}). Sur les
7 fiches sans introduction, soit 69 passages, \code{chunk\_text} ne compte que deux lignes. C'est sur
\code{chunk\_text} qu'est calculée la colonne \code{embeddings\_bge-m3}.

\begin{figure}[htbp]
\centering
\begin{tikzpicture}
  \node[morceau, draw=accent, fill=accent!9, text width=3.2cm] (t) at (0,0)
    {\var{title}\\[2pt]{\color{gris}ligne 1}\\[2pt]« Le bulletin de paie »};
  \node[morceau, draw=accent, fill=accent!9, text width=5.2cm] (i) at (4.25,0)
    {\var{introduction}\\[2pt]{\color{gris}ligne 2}\\[2pt]« Au moment du versement de son salaire, un bulletin de paie doit être remis à chaque salarié… »};
  \node[morceau, draw=accent, fill=accent!9, text width=5.5cm] (x) at (10.55,0)
    {\var{text}\\[2pt]{\color{gris}ligne 3}\\[2pt]« Aucune mention relative à l'exercice du droit de grève… »};
  \node[sf, font=\sffamily\bfseries\fontsize{12}{12}\selectfont, text=accent] at ($(t.east)!0.5!(i.west)$) {+};
  \node[sf, font=\sffamily\bfseries\fontsize{12}{12}\selectfont, text=accent] at ($(i.east)!0.5!(x.west)$) {+};
  \node[morceau, draw=accent2, line width=0.8pt, text width=16.25cm, align=center] (c) at (0,-2.55)
    {\var{chunk\_text} \quad {\color{gris}= les trois lignes réunies, séparées par un retour à la ligne \textbullet{} 619 caractères pour ce passage}};
  \draw[fl] (t.south) -- (t.south |- c.north);
  \draw[fl] (i.south) -- (i.south |- c.north);
  \draw[fl] (x.south) -- (x.south |- c.north);
  \node[sf, anchor=north west, text=rouge, font=\sffamily\fontsize{7}{8.4}\selectfont] at (0,-3.35)
    {La section (\texttt{context}) n'est pas incluse, contrairement à ce qu'annonce la fiche de l'éditeur.};
\end{tikzpicture}
\caption{Composition de \code{chunk\_text}, vérifiée sur les 5\,702 passages, avec l'exemple du passage
n°~10 de la fiche « Le bulletin de paie ».}
\label{fig:composition}
\end{figure}

\subsection{Méthode de découpage appliquée par l'éditeur}\label{sec:decoupage}

D'après la fiche du jeu de données, \code{text} résulte d'un découpage par la fonction de LangChain
\code{Recursive\allowbreak Character\allowbreak Text\allowbreak Splitter}, avec les paramètres suivants.

\begin{tabletype}
\begin{tabularx}{\textwidth}{@{}L{4cm}Y@{}}
\toprule
\thead{Paramètre} & \thead{Valeur} \\
\midrule
\code{chunk\_size} & 1\,024 \\
\code{chunk\_overlap} & 0 \\
\code{length\_function} & tokeniseur de \code{bge-m3} \\
\bottomrule
\end{tabularx}
\end{tabletype}

L'absence de chevauchement signifie que les passages d'une fiche, mis bout à bout dans l'ordre de
\code{chunk\_index}, reconstituent son contenu. MediaTech fournit un tutoriel de reconstitution à cet
effet.

\subsection{Colonnes utiles au projet}

Au regard des sorties exigées par le cas A3 et des besoins du cas B1, les colonnes se répartissent
ainsi.

\begin{tabletype}
\begin{tabularx}{\textwidth}{@{}L{5.2cm}Y@{}}
\toprule
\thead{Usage} & \thead{Colonnes} \\
\midrule
Identification du passage & \code{chunk\_id}, \code{doc\_id}, \code{chunk\_index} \\
Contenu à encoder & \code{text} ou \code{chunk\_text}, choix relevant du domaine D2 \\
Référence de la section d'origine & \code{context} \\
Lien vers le document complet & \code{url}, \code{title} \\
Suivi de version & \code{chunk\_xxh64}, \code{date} \\
Sans usage prévu & \code{surtitle} et \code{source}, constantes ; \code{embeddings\_bge-m3}, modèle différent du nôtre \\
\bottomrule
\end{tabularx}
\end{tabletype}

La colonne \code{embeddings\_bge-m3} représente à elle seule 72,8 millions de caractères, contre
5,3 millions pour la colonne \code{text}. Elle a été calculée avec \code{bge-m3}, que le jalon J2 a écarté
pour le volet A (543,3 Mio, hors budget), et ne peut pas être comparée aux vecteurs de
\code{multilingual-e5-small}.

\subsection{Écarts entre la documentation de l'éditeur et le fichier}

Trois écarts ont été constatés entre la fiche du jeu de données et le fichier réel.

\begin{enumerate}
\item La fiche annonce une colonne \code{surtitre} ; la colonne réelle s'appelle \code{surtitle}.
\item La fiche présente \code{doc\_id} comme l'identifiant du site source sans en préciser le format. Il en
existe trois formes : 257 fiches de la forme \code{article} suivi d'un nombre, 3 identifiants purement
numériques (\code{12399}, \code{12509}, \code{12510}) et 1 identifiant alphanumérique
(\code{article-ACFg1E}).
\item La fiche annonce que \code{chunk\_text} inclut les valeurs de \code{context}. Sur les 5\,702 passages,
\code{chunk\_text} se compose en réalité du titre, de l'introduction et du contenu, sans la section.
\end{enumerate}

Ces écarts sont sans conséquence sur la lecture du fichier, mais tout script doit s'appuyer sur les noms
de colonnes réels, et le choix de la colonne à encoder doit tenir compte de la composition réelle de
\code{chunk\_text}.

% ====================================================================
\section{Filtrage du corpus}\label{sec:filtrage}

\subsection{Principe}

Le filtrage retire du corpus les passages inutilisables : éléments de navigation du site, listes de
liens vers d'autres pages et restes de découpe. Aucune fiche n'est retirée pour être hors périmètre :
les 261 fiches proviennent toutes du site du ministère du Travail et traitent toutes du droit du
travail ou de l'emploi.

Trois principes guident le filtrage.

\begin{itemize}
\item \textbf{Retirer, jamais modifier.} Une règle retire une ligne entière ; aucun texte n'est
corrigé. Le corpus filtré est donc un sous-ensemble exact du fichier d'origine, ce qui le rend simple à
vérifier.
\item \textbf{Évaluer sur \code{text}.} Les règles s'appliquent à la colonne \code{text}, qui contient
le contenu propre du passage. Le titre et l'introduction que \code{chunk\_text} leur ajoute sont communs
à tous les passages d'une fiche : ils ne rendent pas un passage utile. Évaluée sur \code{chunk\_text}, la
règle R3 ne détecterait d'ailleurs plus aucun fragment, puisque le plus court \code{chunk\_text} compte
140 caractères.
\item \textbf{Convenir aux deux encodages possibles.} Le domaine D2 n'a pas encore choisi la colonne à
encoder. Une ligne retirée l'est avec ses deux colonnes, et la règle dont l'intérêt dépend de ce choix
(R4) signale les passages sans les retirer.
\end{itemize}

\subsection{Règles de filtrage}

\textbf{R1 : bouton de navigation seul.} Retire les passages dont le texte se réduit à « Échanger avec le
bon conseiller pour votre entreprise sur Conseillers-Entreprises Service Public », ou à sa variante
« Conseillers Entreprise ». Ce lien de navigation figure au bas de nombreuses fiches et a été extrait
comme un passage à part entière. Il ne répond à aucune question, et ses exemplaires identiques
remonteraient ensemble dans les résultats d'une recherche. \emph{139 passages concernés.}

\textbf{R2 : sections de liens vers d'autres pages.} Retire les passages des sections « Services en
ligne », « À lire sur ce site » et « À lire sur service-public ». Ces sections ne contiennent pas
d'explications mais des renvois : le début du résumé d'autres fiches, coupé par des points de
suspension et suivi de sa rubrique et de sa date, par exemple « La prévention des risques psychosociaux
(RPS) Les risques psychosociaux sont définis comme un risque pour la santé physique et mentale des
travailleurs. Leurs causes sont à\dots{} Risques professionnels Date de mise à jour le 23 juillet 2024 ».
Ces aperçus sont tronqués, et les fiches vers lesquelles ils renvoient figurent en entier dans le corpus.
\emph{343 passages concernés : 171 « Services en ligne », 139 « À lire sur ce site », 33 « À lire sur
service-public ».}

\textbf{R3 : fragments de moins de 50 caractères.} Retire les passages dont le texte compte moins de
50 caractères. Ce sont des restes de découpe, dépourvus de sens isolément : « bis. », « 2026 »,
« de paie », « l'enfant. ». Le seuil résulte d'un compromis : à 30 caractères, des bribes comme
« rémunération au moins équivalente. » seraient conservées ; à 100 caractères, des références
juridiques utiles au volet B, telles que « Articles L. 1223-4, L. 1234-9 à L. 1234-11\dots{} du Code du
travail », seraient retirées. Il est réglable par l'option \code{-{}-seuil} du script.
\emph{76 passages concernés.}

\textbf{R4 : doublons exacts, signalés sans être retirés.} Parmi les passages conservés, marque dans la
nouvelle colonne \code{doublon\_text} ceux dont le texte est identique à celui d'un passage précédent,
dans l'ordre du fichier ; la première occurrence n'est pas marquée. Ces doublons se trouvent tous dans
des fiches différentes, et sont surtout des aperçus de liens comme « Glossaire des mots et des
institutions Détachement international | Date de mise à jour le 27 juillet 2020 ». Leur intérêt dépend
de la colonne encodée : encodés à partir de \code{text}, ils produisent des vecteurs identiques qui
occupent plusieurs places d'un même classement ; encodés à partir de \code{chunk\_text}, ils restent
distincts, puisque le titre de leur fiche diffère. Le domaine D2 les écarte ou non selon son choix ;
l'option \code{-{}-sans-doublons} du script les retire directement, et les ajoute alors, sous la règle R4,
à la liste des passages retirés.
\emph{40 passages marqués.}

\subsection{Ce qui est volontairement conservé}

Deux types de sections, proches en apparence des listes de liens, ne sont pas visés par R2.

\begin{itemize}
\item \textbf{« Textes de référence »} : références d'articles du Code du travail et de textes
réglementaires. Elles sont utiles au volet B, qui doit citer ses sources. 181 passages sont conservés.
\item \textbf{« Qui contacter »} : organismes et services à contacter, qui répondent aux questions du
type « à qui m'adresser ? ». 49 passages sont conservés.
\end{itemize}

Les passages de ces sections qui comptent moins de 50 caractères sont néanmoins retirés par R3.

\subsection{Nombre de passages conservés et volumétrie}

Un passage peut relever de plusieurs règles : 136 boutons de navigation se trouvent dans la section
« Services en ligne » (R1 et R2), et 4 fragments courts dans une section de liens (R2 et R3). Chaque
passage retiré est compté une seule fois, sous la première règle qui le concerne dans l'ordre R1, R2, R3.

\begin{tabletype}
\begin{tabularx}{\textwidth}{@{}Y L{2.6cm}L{2.9cm}@{}}
\toprule
\thead{Règle} & \thead{Concernés, pris isolément} & \thead{Retirés à ce titre} \\
\midrule
R1, bouton de navigation seul & 139 & 139 \\
R2, sections de liens & 343 & 207 \\
R3, fragments de moins de 50 caractères & 76 & 72 \\
\midrule
Total retiré & & \textbf{418} (7,3 \%) \\
\bottomrule
\end{tabularx}
\end{tabletype}

La volumétrie est mesurée par trois indicateurs : le nombre de passages, le nombre de fiches
distinctes, et le nombre total de caractères de la colonne \code{text}, qui représente la quantité de
contenu exploitable.

\begin{tabletype}
\begin{tabularx}{\textwidth}{@{}Y L{2cm}L{1.8cm}L{3.6cm}@{}}
\toprule
\thead{Corpus} & \thead{Passages} & \thead{Fiches} & \thead{Caractères de \code{text}} \\
\midrule
Brut, avant filtrage & 5\,702 & 261 & 5\,295\,169 \\
Filtré par R1, R2 et R3 & \textbf{5\,284} & \textbf{261} & 5\,159\,728 (−2,6 \%) \\
Filtré, doublons marqués par R4 écartés & 5\,244 & 261 & 5\,136\,460 (−3,0 \%) \\
\bottomrule
\end{tabularx}
\end{tabletype}

Le filtrage retire 7,3 \% des passages mais seulement 2,6 \% du texte : ce qu'il enlève est presque
entièrement constitué de passages très courts ou sans contenu propre. Aucune fiche n'est perdue.

Le corpus filtré conserve 13 des 14 colonnes d'origine : la colonne \code{embeddings\_bge-m3}, calculée
avec un autre encodeur que celui du projet, n'est pas reprise, ce qui ramène le fichier de 25,9 Mio à
2,7 Mio. S'y ajoutent trois colonnes :
\code{doublon\_text} (règle R4), et \code{tokens\_text} et \code{tokens\_chunk\_text}, qui donnent le nombre
de tokens de chaque forme du passage (section~\ref{sec:troplong}).
Les 418 passages retirés sont listés, avec la règle appliquée et son motif, dans le fichier
\code{data/processed/\allowbreak passages\_retires.csv}.

% ====================================================================
\section{Contrôle manuel d'un échantillon}\label{sec:controle}

\subsection{Méthode}

Deux échantillons ont été tirés au hasard par le script \code{scripts/echantillon\_controle.py}, avec
une graine fixe (27) qui rend le tirage reproductible :

\begin{itemize}
\item \textbf{échantillon A} : 50 passages conservés par le filtrage, issus de 40 fiches différentes,
pour vérifier que le filtrage n'a pas laissé passer de passages inutilisables ;
\item \textbf{échantillon B} : 20 passages retirés, répartis entre les règles (7 au titre de R1, 7 de R2,
6 de R3), pour vérifier que chaque retrait était justifié.
\end{itemize}

La colonne à encoder n'étant pas encore choisie, chaque passage a été examiné sous ses deux formes,
\code{text} et \code{chunk\_text}, avec un verdict pour chacune. Le verdict porte sur l'utilité du
passage : un passage coupé en début ou en fin mais dont le contenu reste exploitable est jugé conforme,
la coupure étant relevée séparément dans la colonne \code{coupure} (section~\ref{sec:coupures}). Le
verdict de \code{chunk\_text} reprend les défauts de \code{text}, que le titre et l'introduction ne
corrigent pas, et y ajoute les défauts propres à cette forme, en particulier une longueur supérieure à
512 tokens, au-delà de laquelle l'encodeur tronquerait le passage (section~\ref{sec:troplong}). Le
nombre de tokens de chaque passage est compté par \code{scripts/prepare\_corpus.py}.

Le fichier \code{controle\_echantillon.xlsx}, rangé dans \code{docs/corpus/}, contient la grille
complète, avec un commentaire pour chaque passage. Le contrôle a été réalisé avec l'assistance
d'un outil d'intelligence artificielle, puis relu et validé par Vaneck DAGAR.

\subsection{Résultats}

\begin{tabletype}
\begin{tabularx}{\textwidth}{@{}Y L{2.6cm}L{2.6cm}@{}}
\toprule
\thead{Échantillon A, passages conservés} & \thead{\code{text}} & \thead{\code{chunk\_text}} \\
\midrule
Passages conformes & \textbf{35 sur 50 (70 \%)} & \textbf{30 sur 50 (60 \%)} \\
Section d'origine erronée & 8 & 8 \\
Bruit de navigation ou liste de liens & 5 & 5 \\
Trop long pour l'encodeur & 0 & 6 \\
Fragment incompréhensible & 1 & 0 \\
Tableau aplati & 1 & 1 \\
\midrule
Passages coupés en début ou en fin & 46 sur 50 & 46 sur 50 \\
\bottomrule
\end{tabularx}
\end{tabletype}

Deux passages de \code{chunk\_text} jugés non conformes pour une autre raison dépassent aussi la
longueur limite : 8 passages sur 50 comptent plus de 512 tokens. Un neuvième en compte exactement 512 ;
il tient dans la limite, mais serait tronqué de 2 tokens si l'on ajoutait le préfixe \code{passage:}
attendu par les modèles E5.

\textbf{Échantillon B.} 19 retraits sur 20 sont justifiés, pour les deux formes du passage. Le seul
retrait injustifié concerne un passage de 32 caractères de la section « Qui contacter ? »
(« France Travail · Votre DREETS-DDETS »), retiré par R3 alors qu'il indique utilement à qui
s'adresser. C'est une limite du seuil de 50 caractères (section~\ref{sec:limites}).

\subsection{Interprétation}

Les sections erronées dépendent de l'extraction de l'éditeur et non de la colonne encodée. Hors de ce
défaut commun, 7 passages sur 50 (14 \%) posent problème sous la forme \code{text}, contre 13 sur 50
(26 \%) sous la forme \code{chunk\_text}. L'écart tient surtout à la longueur : \code{chunk\_text}
répète l'introduction de la fiche, qui compte jusqu'à 2\,056 caractères. L'avantage attendu de
\code{chunk\_text}, donner un contexte aux passages qui en manquent, n'a été décisif que pour un passage
sur 50.

Ces proportions sont à lire comme des ordres de grandeur : sur 50 passages, une proportion de 70 \%
est connue à environ 13 points près, avec un niveau de confiance de 95 \%.

Le contrôle est donc favorable à l'encodage de \code{text}. Une composition intermédiaire, formée du
titre de la fiche suivi du texte du passage, conserverait un contexte sans dépasser la limite de
l'encodeur : elle compte 438 tokens au plus. Le choix revient au domaine D2.

% ====================================================================
\section{Anomalies et limites}\label{sec:anomalies}

Cette section recense les défauts observés dans le corpus. Les mesures portent sur le corpus filtré
décrit en section~\ref{sec:filtrage}, soit 5\,284 passages, sauf mention contraire.

\subsection{Passages coupés en milieu de phrase}\label{sec:coupures}

\textbf{Constat.} La plupart des passages ne commencent pas au début d'une phrase ou ne se terminent pas
à sa fin. Le passage \code{12509\_21} de la fiche « La période de reconversion » en donne un exemple :
il commence par « que se passe-t-il ? La non-réponse dans le délai vaut décision de refus de prise en
charge », suite d'une question posée dans le passage précédent, et se termine par « Néanmoins, ».

\begin{tabletype}
\begin{tabularx}{\textwidth}{@{}Y L{2.6cm}L{2.2cm}@{}}
\toprule
\thead{Indicateur} & \thead{Passages} & \thead{Part} \\
\midrule
Début en milieu de phrase (première lettre minuscule) & 3\,113 & 58,9 \% \\
Fin sans ponctuation forte (point, point d'interrogation ou d'exclamation, deux-points, point-virgule, guillemet ou parenthèse fermants) & 3\,896 & 73,7 \% \\
Début ou fin coupé & 4\,537 & 85,9 \% \\
\midrule
Pour mémoire, début en milieu de phrase dans le corpus brut & 3\,167 sur 5\,702 & 55,5 \% \\
Pour mémoire, début ou fin coupé dans l'échantillon contrôlé & 46 sur 50 & 92 \% \\
\bottomrule
\end{tabularx}
\end{tabletype}

Ces indicateurs sont approximatifs. Un passage peut commencer par une majuscule tout en poursuivant une
phrase entamée dans le passage précédent, par exemple lorsqu'il débute par un sigle : la part réelle des
coupures est donc au moins égale aux valeurs indiquées.

\textbf{Cause.} Le défaut provient du découpage appliqué par l'éditeur (section~\ref{sec:decoupage}). Lorsqu'une
section de fiche dépasse la taille maximale d'un passage, elle est coupée en tranches successives sans
égard aux limites des phrases. Les passages concernés ne sont donc pas du bruit : ce sont les suites des
passages précédents, et ils portent l'essentiel du contenu juridique.

\textbf{Pourquoi ne pas les retirer.} Retirer ces passages viderait le corpus de la majeure partie de son
contenu, comme le montre le tableau suivant.

\begin{tabletype}
\begin{tabularx}{\textwidth}{@{}Y L{2.6cm}L{2.2cm}L{2.4cm}@{}}
\toprule
\thead{Passages retirés} & \thead{Passages restants} & \thead{Fiches restantes} & \thead{Texte conservé} \\
\midrule
Début en milieu de phrase & 2\,171 (41 \%) & 261 & 37 \% \\
Début ou fin coupé & 747 (14 \%) & 242 & 9 \% \\
\bottomrule
\end{tabularx}
\end{tabletype}

Ces passages sont donc conservés, et la coupure est signalée pour chaque passage de l'échantillon de
contrôle dans la colonne \code{coupure} de \code{docs/corpus/controle\_echantillon.xlsx}.

\subsubsection{Effet sur l'encodage}

La coupure n'empêche pas l'encodage, mais elle en dégrade la qualité sur certains points. Ses effets
sont de nature et de gravité différentes.

\begin{tabletype}
\begin{tabularx}{\textwidth}{@{}L{4.3cm}Y L{2.3cm}@{}}
\toprule
\thead{Effet} & \thead{Explication} & \thead{Gravité} \\
\midrule
Fragment de phrase en tête de passage & L'encodeur résume l'ensemble du passage en faisant la moyenne des représentations de ses tokens. Quelques mots isolés en début de passage pèsent peu face aux quelque deux cents autres : le sens global du vecteur est préservé. & faible \\
Idée partagée entre deux passages & Une règle commencée à la fin d'un passage et achevée au début du suivant n'est complètement représentée par aucun des deux vecteurs. Une question portant précisément sur cette règle risque d'être mal rapprochée de l'un comme de l'autre. & moyenne \\
Sujet absent du passage & Un passage qui commence par « Cette transaction… » ou « à les mettre en concurrence » a perdu ce dont il parle. Son vecteur représente la notion en général, sans son contexte, et se rapproche moins bien d'une question formulée avec ce contexte. La colonne \code{chunk\_text} compense en partie ce défaut, puisqu'elle ajoute le titre et l'introduction de la fiche. & moyenne \\
\bottomrule
\end{tabularx}
\end{tabletype}

\textbf{Conséquences au-delà de l'encodage.} Dans le volet A, l'utilisateur voit s'afficher un
résultat qui commence en milieu de phrase, sans savoir à quoi il se rapporte. Dans le volet B, le
modèle de langage reçoit un passage privé de son début : il peut citer la bonne source tout en se
trompant sur ce à quoi la règle s'applique.

\textbf{Conclusion.} La coupure n'est pas bloquante pour l'encodage, mais elle devrait réduire le rappel
sur les questions précises. Son effet peut être mesuré : il suffit de comparer le Recall@k obtenu sur le
corpus actuel à celui d'un corpus re-segmenté, avec le même jeu de questions.

\textbf{Piste de correction.} Le découpage de l'éditeur étant sans chevauchement, chaque section peut
être reconstituée en recollant ses passages dans l'ordre de \code{chunk\_index}, puis redécoupée en
respectant les fins de phrase et la limite de 512 tokens de \code{multilingual-e5-small}. Ce travail
relève du domaine D1 mais sort du périmètre du ticket \#27, qui porte sur le filtrage. Il modifie le
nombre de passages et concerne le domaine D2 : il doit être proposé en réunion hebdomadaire et faire
l'objet d'un ticket dédié avant d'être engagé.

\subsection{Sections d'origine mal attribuées}

\textbf{Constat.} Vingt-cinq fiches comportent une section consacrée à une vidéo, par exemple
« Le contrat à durée déterminée (CDD) en vidéo (web série droit du travail) ». Ces sections regroupent
793 passages (15,0 \%). Seuls 79 d'entre eux portent une marque de transcription orale (« Bonjour »,
« épisode », « retranscription », « cette vidéo ») ; les 714 autres (13,5 \%) sont des textes juridiques
écrits, sans rapport avec la vidéo. L'extraction de l'éditeur semble avoir rangé la suite de la page sous
le dernier intertitre rencontré. L'échantillon contrôlé confirme ce défaut pour 8 passages sur 50.
D'autres passages, à cheval sur plusieurs sous-sections, ne portent que l'intitulé de la première.

\textbf{Conséquences.} La colonne \code{context} n'est pas fiable. Ses effets diffèrent selon l'usage
qui en est fait.

\begin{tabletype}
\begin{tabularx}{\textwidth}{@{}L{3.6cm}Y@{}}
\toprule
\thead{Usage} & \thead{Effet d'une section erronée} \\
\midrule
Encodage et recherche & Aucun : \code{context} ne figure ni dans \code{text} ni dans \code{chunk\_text}, et un passage mal attribué est retrouvé comme s'il était bien attribué. L'effet deviendrait réel si le domaine D2 choisissait d'encoder une composition incluant la section. \\
Volet A, affichage des résultats & Le cas A3 affiche avec chaque résultat la référence de sa section d'origine. Elle serait fausse pour environ 13,5 \% des passages, ce qui induirait l'utilisateur en erreur. \\
Volet B, citations & L'assistant sourcé peut citer une section qui ne correspond pas au passage. Si la section est fournie au modèle de langage avec chaque passage, une étiquette erronée peut aussi l'orienter à tort. \\
Volet B, évaluation & La métrique « taux de citations valides » serait faussée si la validité d'une citation était jugée sur la section. Elle doit l'être sur l'identifiant du passage (\code{chunk\_id}) et l'adresse de la fiche (\code{url}). \\
\bottomrule
\end{tabularx}
\end{tabletype}

Le titre (\code{title}) et l'adresse (\code{url}) de la fiche restent toujours justes : ils sont définis au
niveau de la fiche et ne dépendent pas de l'extraction des sections. Ils constituent la référence fiable
à afficher et à citer. La section peut être présentée comme une information indicative, ou masquée pour
les intitulés de vidéo qui ne portent aucune marque de transcription.

\subsection{Bruit de navigation restant après filtrage}

\textbf{Constat.} 327 passages conservés (6,2 \%) contiennent encore des éléments de navigation :
183 le bouton « Échanger avec le bon conseiller », 223 des aperçus de liens terminés par « Date de mise à
jour ». Les règles R1 et R2 ne les ont pas détectés, pour deux raisons : le bouton y est mêlé à d'autres
textes, alors que R1 ne retire que les passages qui s'y réduisent ; et ces passages sont rangés sous une
section mal attribuée, que R2 ne reconnaît pas comme une liste de liens. Des listes de documents
téléchargeables, rangées sous « Pour aller plus loin », relèvent du même cas. L'échantillon contrôlé en
compte 5 sur 50.

\textbf{Effet sur l'encodage.} Le bruit dilue le vecteur du passage : une partie de sa représentation
porte sur des titres d'autres fiches plutôt que sur son propre contenu. Le passage peut alors remonter
pour des questions qui relèvent d'autres fiches.

\subsection{Longueur excessive de chunk\_text}\label{sec:troplong}

\textbf{Mesure.} Le script \code{scripts/prepare\_corpus.py} compte les tokens de chaque passage avec le
tokeniseur de l'encodeur retenu au jalon J2, \code{Xenova/multilingual-e5-small}, figé sur un commit.
Le décompte inclut les deux tokens spéciaux de début et de fin ; il est enregistré dans les colonnes
\code{tokens\_text} et \code{tokens\_chunk\_text} du corpus filtré.

\begin{tabletype}
\begin{tabularx}{\textwidth}{@{}Y L{1.8cm}L{1.8cm}L{3.9cm}@{}}
\toprule
\thead{Forme du passage} & \thead{Médiane} & \thead{Maximum} & \thead{Au-delà de 512 tokens} \\
\midrule
\code{text} & 281 & 431 & 0 \\
\code{chunk\_text} & 394 & 804 & \textbf{619 (11,7 \%)} \\
\code{chunk\_text} avec le préfixe \code{passage:} & & & 637 (12,1 \%) \\
Titre de la fiche et \code{text} & 292 & 438 & 0 \\
\bottomrule
\end{tabularx}
\end{tabletype}

L'écart vient de l'introduction de la fiche, répétée dans chaque \code{chunk\_text} : elle dépasse
1\,000 caractères pour 21 fiches, et atteint 2\,056 caractères. Les modèles E5 attendent que chaque
passage soit précédé du préfixe \code{passage:}, qui ajoute 2 tokens : pris seul, il en compte 3, mais
collé au texte, son dernier token fusionne avec le premier mot. La troisième ligne du tableau en tient
compte ; le décompte est fait directement sur le texte préfixé.

\textbf{Effet sur l'encodage.} L'encodeur \code{multilingual-e5-small} ne traite que les 512 premiers
tokens d'une entrée et ignore la suite. Le titre et l'introduction venant en tête de \code{chunk\_text},
c'est la fin du passage, c'est-à-dire son contenu propre, qui serait perdue. Pour les 619 passages
concernés, la troncature supprime 55 tokens en médiane, et jusqu'à 292, soit en médiane 19 \% du
contenu propre du passage.

La première mesure de ce défaut reposait sur une estimation de quatre caractères par token, qui
désignait 849 passages de plus de 2\,000 caractères. Le décompte exact la corrige : l'estimation
surévaluait de 37 \% le nombre de passages concernés.

\subsection{Défauts de mise en forme du texte}

\begin{tabletype}
\begin{tabularx}{\textwidth}{@{}L{4.4cm}Y L{2.6cm}@{}}
\toprule
\thead{Défaut} & \thead{Exemple} & \thead{Passages} \\
\midrule
Intertitre collé au texte & « prise en charge.Systèmes d'informationUn dépôt du CERFA » & environ 3\,439 (65,1 \%) \\
Adresse web coupée par une espace & « ameli. fr », « epide. fr » & 146 (2,8 \%) \\
Transcription orale d'une vidéo & « Bonjour, c'est Honorine et aujourd'hui je vais vous présenter\dots{} » & 79 (1,5 \%) \\
Espace avant un point ou une virgule & « un bulletin de paie doit être remis à chaque salarié . » & 30 (0,6 \%) \\
Entité HTML non décodée & « d\&\#8217;un congé » au lieu de « d'un congé » & 5 \\
Tableau aplati & « CVECDIPublic viséOuvert aux demandeurs d'emploi\dots{} » & 2 sur 50 dans l'échantillon \\
Texte de lien dans l'introduction & « (Lien vers « La définition du CSE et cadre de mise en place ») » & 2 fiches \\
\bottomrule
\end{tabularx}
\end{tabletype}

Les mesures reposent sur des expressions régulières et sont approximatives, en particulier pour les
intertitres collés. Les tableaux aplatis n'ont pas pu être mesurés sur l'ensemble du corpus faute
d'indicateur fiable.

\textbf{Effet sur l'encodage.} La plupart de ces défauts sont mineurs : le tokeniseur découpe les mots
en sous-unités et tolère une espace mal placée. Les intertitres collés n'ont qu'un effet négligeable sur
l'encodeur, qui découpe « informationUn » en « information » puis « Un », comme il le ferait avec une
espace ; ils empêchent en revanche de séparer un intertitre de la phrase qui le suit lors d'un découpage
en phrases. Les tableaux aplatis sont plus gênants : l'association entre les cellules et leurs colonnes
est perdue, et le passage devient ambigu.

\subsection{Structure et contenu du corpus}

\begin{itemize}
\item \textbf{Déséquilibre entre fiches.} Une fiche compte de 1 à 196 passages conservés, 13 en médiane.
Les dix fiches les plus longues représentent 20,2 \% du corpus : les sujets qu'elles traitent, comme
l'activité partielle de longue durée, seront surreprésentés dans les résultats de recherche.
\item \textbf{Doublons entre fiches.} 40 passages reprennent à l'identique le texte d'un passage d'une
autre fiche ; ils sont signalés par R4 (section~\ref{sec:filtrage}).
\item \textbf{Fiches sans introduction.} 7 fiches, soit 65 passages conservés, n'ont pas d'introduction.
Le \code{chunk\_text} de leurs passages ne compte alors que deux lignes et apporte moins de contexte.
\item \textbf{Contenus d'actualité inégale.} Les fiches ont été publiées ou mises à jour entre le
17 novembre 2008 et le 27 août 2026 ; 34 datent d'avant 2020. Certaines conservent des éléments
historiques, comme les constats de 2008 dans la fiche du plan santé au travail 2026-2030.
\end{itemize}

\subsection{Limites}\label{sec:limites}

\textbf{Limites du corpus.}

\begin{itemize}
\item Le corpus provient d'une seule source, le site du ministère du Travail : il ne couvre ni les
conventions collectives, ni la jurisprudence au-delà de ce que les fiches en citent.
\item Il reflète l'état des fiches au 11 septembre 2026. Les évolutions réglementaires postérieures, et
les mises à jour publiées depuis par l'éditeur, ne sont pas prises en compte.
\item La colonne \code{context} n'est pas fiable, et la colonne \code{embeddings\_bge-m3} n'est pas
utilisable par le projet, puisqu'elle a été calculée avec un autre encodeur.
\end{itemize}

\textbf{Limites du filtrage.}

\begin{itemize}
\item Les règles R1 et R2 s'appuient sur le texte exact du bouton et sur l'intitulé des sections : elles
laissent passer le bruit mêlé à d'autres textes ou rangé sous une section mal attribuée, soit 6,2 \% des
passages conservés.
\item Le seuil de 50 caractères de R3 retire parfois une information courte mais utile, comme un
organisme à contacter ou la référence d'un décret : un retrait injustifié sur 20 dans l'échantillon
contrôlé.
\item Le filtrage ne corrige pas les coupures en milieu de phrase, qui relèvent d'une re-segmentation
hors du périmètre du ticket \#27.
\end{itemize}

\textbf{Limites des mesures.}

\begin{itemize}
\item Les indicateurs d'anomalie reposent sur des expressions régulières : ce sont des estimations,
qui peuvent comporter des faux positifs et des faux négatifs.
\item Les longueurs en tokens sont mesurées avec le tokeniseur du modèle d'origine. Si le vocabulaire
de l'encodeur est réduit, comme le prévoit le jalon J2, le découpage en tokens changera et le décompte
devra être refait.
\item L'échantillon contrôlé compte 70 passages : ses proportions sont des ordres de grandeur, à environ
13 points près pour l'échantillon A.
\end{itemize}

\end{document}
